\documentclass[aspectratio=169, 10pt]{beamer}

\usepackage[utf8]{inputenc}
\usepackage[T1]{fontenc}
\usepackage{xcolor}
\usepackage{graphicx}

\IfFileExists{opensans.sty}{\usepackage[default,scale=0.95]{opensans}}{}

\definecolor{amritaMaroon}{HTML}{A4123F}
\definecolor{macroBlue}{HTML}{2B5B84}
\definecolor{microOrange}{HTML}{D97925}

\usetheme{Madrid}
\usecolortheme[named=amritaMaroon]{structure}
\setbeamertemplate{navigation symbols}{}
\setbeamertemplate{footline}[frame number]

\begin{document}

% --- SLIDE 1 ---
\begin{frame}{Cooperative Co-Evolution: Splitting the Brain}

\textbf{The Problem:} Searching a massive 20-dimensional continuous space is computationally expensive. Good behaviors are often lost if another part of the brain mutates poorly (a "tragic flaw").

\vspace{0.3cm}
\textbf{The Solution:} Decompose the single "monolithic" brain into two distinct, specialized sub-populations. 

\vspace{0.4cm}
\begin{columns}[T]
    \begin{column}{0.48\textwidth}
        \begin{block}{\textbf{\textcolor{macroBlue}{Population A: Macro / Economy}}}
            \begin{itemize}
                \item Evolves \textbf{8 parameters} independently.
                \item Focuses on base building, worker allocation, and resource timing.
                \item \textit{e.g., worker\_ratio, barracks\_timing}
            \end{itemize}
        \end{block}
    \end{column}
    
    \begin{column}{0.48\textwidth}
        \begin{block}{\textbf{\textcolor{microOrange}{Population B: Micro / Combat}}}
            \begin{itemize}
                \item Evolves \textbf{12 parameters} independently.
                \item Focuses on unit targeting, aggression, and tactical combat behavior.
                \item \textit{e.g., target\_priority, focus\_fire}
            \end{itemize}
        \end{block}
    \end{column}
\end{columns}

\vspace{0.4cm}
\textit{Why? Mathematically, searching two smaller spaces (8D + 12D) converges exponentially faster than searching one large space (20D).}

\end{frame}

% --- SLIDE 2 ---
\begin{frame}{The Cooperative "Handshake" (Evaluation Phase)}

Since the game engine strictly requires all 20 weights to play a match, the sub-populations must \textbf{cooperate} to be evaluated.

\vspace{0.4cm}

\textbf{\textcolor{macroBlue}{1. Evaluating an Economy Agent (Pop A):}}
\begin{itemize}
    \item We take an 8-weight Macro agent from Pop A.
    \item We \textbf{borrow} the 12-weight Micro parameters from the \textbf{current best} agent in Pop B.
    \item We combine them (8 + 12 = 20 weights) and play the game. The win/loss fitness is assigned \textbf{only} to the Macro agent.
\end{itemize}

\vspace{0.2cm}
\textbf{\textcolor{microOrange}{2. Evaluating a Combat Agent (Pop B):}}
\begin{itemize}
    \item We take a 12-weight Micro agent from Pop B.
    \item We \textbf{borrow} the 8-weight Macro parameters from the \textbf{current best} agent in Pop A.
    \item We combine them and play the game. The fitness is assigned \textbf{only} to the Micro agent.
\end{itemize}

\vspace{0.3cm}
\begin{alertblock}{The Advantage}
This architecture ensures that a brilliant economic strategy isn't penalized just because it was randomly generated with a terrible combat behavior. Good partial solutions are protected and combined!
\end{alertblock}

\end{frame}

\end{document}
